\section{Benchmark}

To evaluate the security auditing capabilities of existing MCP vulnerability detection tools, we construct the \bench benchmark. We present a detailed description of the dataset construction process and its statistics.


\subsection{Data Collection}


% 1. 本基准测试集共包含 270 条经过人工校验的测试样本。

% 2. 在编程语言分布上，数据集涵盖了当前 MCP 生态中最主要的两门语言：Python 与TypeScript 。选择这两门语言作为基准，是因为 Python 与 TypeScript 是目前 MCP 开发的两大主流语言，占据了 95\% 以上的生态份额。 ----提供reference

XXX. \cite{tensorrt}

% 3. 具体来对于正样本，我们 XXX，最终得到 Python样本XX条，TypeScript样本XXX条。
% 4. 对于负样本(对照数据)，我们 XXX，最终得到 Python样本XX条，TypeScript样本XXX条。

% 170 + 100。

% 1. Python XX大漏洞，常用表现形式

% 2. RCE 25 对 ， 其他 15， 

% 25对 5组，每组间关键触发点不同，每组内，关键触发点相同，但是XX不同。

% Gemini Prompt. 提供prompt, 人工验证

% 85 + 50 
% 85 + 50。

        % 2.4 对照数据
        % 2.4.1 对照数据构建的必要性
        % 构建对照数据的主要目的在于消除LLM在大规模预训练中形成的模式匹配偏见。大语言模型往往对部分关键词（如 os.system、eval、exec）表现出过度敏感，倾向于在不分析上下文过滤逻辑的情况下直接给出漏洞判定。这种行为导致了高虚警率。引入对照数据可以强制模型必须识别出有效的安全防御逻辑（如参数化、白名单校验）与无效防御逻辑（如非递归替换、黑名单）之间的本质区别。此外，对照数据是计算“同答比”指标的基础，只有当模型能够同时在同一功能逻辑下准确区分漏洞版与安全版时，才认为其具备了真正的安全语义理解能力。
        % 2.4.2 对照数据的构建思路
        % 本数据集遵循最小扰动原则进行对照数据的开发，在代码结构上，安全对照版与对应的漏洞版保持极高的相似性。除了核心的脆弱点执行路径或输入清洗逻辑外，其余框架初始化代码、工具定义及非关键代码完全一致。这种设计确保了模型无法通过代码长度、函数名或变量名等统计特征进行猜测，必须聚焦于代码层面的具体逻辑的差异。同时，在修复逻辑上，本文采用了工业界公认的最佳实践。对于命令注入风险，采用参数数组化调用或强制转义；对于 SQL 注入，采用数据库驱动层面的参数化查询；对于路径穿越，采用路径规范化比对及基础目录锚定。这种设计不仅提供了评估基准，还通过起到了代码的示范作用。
        % 2.4.3 对照数据构建过程
        % 构建过程分为漏洞点识别、防御逻辑重构以及人工校验阶段。首先，深入分析 135 条漏洞样本，确定其触发风险的关键节点。随后，针对 Python 与 TypeScript 的语言特性，编写对应的防御逻辑，并确保修复后的代码在 MCP 协议下依然保持工具的可调用性与功能完整性。最后，所有生成的安全样本经过交叉人工审核，确保不含逻辑残留漏洞，并将其 id 设置为 MCP-SAFE- 前缀，实现与漏洞样本的 1:1 逻辑关联。




\begin{table}[!t]
\centering
\small
    \begin{tabular}{m{2.5cm}m{1.5cm}m{2cm}m{1cm}}
\toprule
\textbf{Dataset Identifier} & \textbf{Vulnerability Type} & \textbf{Vulnerability Description} & \textbf{Count}\\
\hline
MCP-SSRF-001-cpl to 005-cpl & SSRF & XXX & XXX \\ 

\midrule
MCP-XXX to MCP-XXX & - & - & 270 \\ 


\midrule
\bottomrule
\end{tabular}
\caption{\bench}
\label{tab:dataset}
\end{table}


\subsection{Dataset Schema}

% 2.2.1 标识与评估字段：id 与 score
% id 字段不仅是数据的唯一索引，更承载了实验分类的功能。通过在 id 中加入前缀（如 MCP-SAFE-）或后缀（如 -cpl），实验程序可以在无需解析代码的情况下快速识别样本的属性，从而实现“同答比”和“复杂比”等指标的自动统计。例如，MCP-PT-009-cpl就代表第009号漏洞类型为PT的漏洞数据，且这条数据为复杂漏洞，而MCP-SAFE-PT-009即为其对照数据
% score 字段用于设置样本分数。本文设定简单漏洞样本识别正确得 1分，复杂漏洞样本识别正确得2分，安全样本识别正确得 1 分。这种差异化赋分策略模拟了真实的审查逻辑：漏报一个关键漏洞导致的系统风险远高于误报一个安全片段。因此，该字段确保了评估分值能更真实地反映被测模型的实际审计能力
% 2.2.2 上下文与标签字段：metadata 与 ground_truth
% metadata 字段提供了样本的背景信息，包括编程语言、漏洞类型的文本描述等。
% ground_truth 字段是评估的客观准则。除布尔类型的漏洞判定外，它还包含了细化的风险分类、对应的 CWE 编号以及漏洞在代码中的具体位置描述。这种结构化的标签系统使得基准测试不仅能判断被测程序 “答得对不对”，还能进一步通过更复杂的解析流程判断被测程序是否“找得准”。
% 2.2.3 assets字段
% assets 字段是本基准测试的核心。该字段通过嵌套数组形式存储了完整的源代码文件。本文构建的 assets 数据在结构上完美契合了真实的 MCP 交互场景。
% 首先是协议框架完整性。 数据集样本并非零散的逻辑代码，而是完整的 MCP 服务端实现。 以 Python 样本（如 MCP-RCE-001-cpl）为例，代码中集成了 mcp.server.fastmcp 框架，并包含了基于 starlette 的异步传输层（SSE）配置逻辑。这种设计模拟了 MCP 服务端在真实环境中的标准部署模式，要求 LLM 在识别漏洞的同时，必须理解协议栈的运行机制。 在 TypeScript 样本（如 MCP-RCE-001-cpl）中，代码严格遵循官方 @modelcontextprotocol/sdk 开发规范，通过 McpServer 类实例化服务端对象，并配置了标准的 StdioServerTransport 传输层。这种对官方 SDK 的原生支持，确保了测试数据与真实生态环境的无缝衔接。
% 其次是工具调用范式还原。 MCP 的核心在于“工具”的定义。assets 字段中的代码通过标准的协议范式封装逻辑。 在 Python 数据中，利用 @mcp.tool() 装饰器将普通函数注册为 MCP 工具，这使得输入参数（如 hostname）直接对应协议层接收到的外部指令。 在 TypeScript 数据中，使用 server.tool() 并配合 zod 库进行模式校验，复现了真实场景下对于参数类型的强约束检查。这种封装方式复现了 LLM 作为代理调用工具时所面临的真实数据流——从工具参数接收用户输入到执行高危操作的全过程。
% 并且每段代码都有闭环运行能力， 每个 assets 字段提供的代码均包含完整的导入依赖说明和主函数入口。 例如，Python 样本中包含了 uvicorn.run() 的启动配置，TypeScript 样本中包含了 server.connect() 的连接逻辑。这意味着测试代码在语法层面是完整且可运行的。这种高度的现实契合性保证了 LLM 在评估过程中处理的是具有真实上下文背景的业务代码


\subsection{Dataset Characteristics}

\textbf{RCE.}


\textbf{SQL Injection.}

\textbf{DoS.}


\subsection{Metrics}



% 1.	基础分类指标（准确率、精确率、召回率）：
% 准确率（Accuracy）衡量模型整体判断的正确比例性。精确率（Precision）侧重于评估模型报警的质量，即在模型判定的漏洞中，真实漏洞所占的比例，用于衡量虚警风险。召回率（Recall）则反映了模型发现潜在漏洞的能力，即在所有真实漏洞中，被成功检出的比例。这三个指标构成了评估模型基础审计能力的基石。
% 2.	同答比 (PCR)： 同答比是衡量模型是否具备真实安全逻辑理解力的核心指标。其计算逻辑是：将数据集按照逻辑关联进行配对（每一对包含一个漏洞样本及其对应的安全样本），只有当模型能够同时准确识别出该组内的漏洞与安全性时，才判定该组“同答正确”。 
% 该指标的作用在于剔除那些依靠关键词匹配或猜测而获得的偶然正确得分，强制模型证明其能够区分微小的防御逻辑差异。







% 3.	复杂比 (CR)： 复杂比用于评估模型在应对诱导性代码时的鲁棒性。其定义为：在模型判定为漏洞的样本集合中，带有 -cpl 后缀的高复杂度样本所占的比例。 
% 高复杂比意味着模型能够有效识别无效过滤、异常回退等逻辑陷阱，是区分初级安全检查与深度安全审计能力的关键指标。
